\begin{resumo}
A segurança do paciente e o rastreamento retrospectivo de eventos adversos em instituições de saúde constituem competências fundamentais e de alta complexidade na formação clínica contemporânea. A metodologia \textit{Global Trigger Tool} (GTT), padronizada pelo \textit{Institute for Healthcare Improvement} (IHI), preconiza a auditoria sistemática de prontuários em um intervalo estrito de vinte minutos por prontuário, combinando a busca de gatilhos com a posterior investigação causal por meio de ferramentas da qualidade. Todavia, a adoção de métodos tradicionais baseados em papel ou planilhas genéricas gera elevada sobrecarga cognitiva, ansiedade temporal e fadiga visual em discentes e docentes. Neste contexto, o SIGEA-GTT (Sistema Inteligente de Gestão de Eventos Adversos -- Global Trigger Tool) foi concebido sob a perspectiva do Design Centrado no Usuário (DCU), em consonância com a norma ISO 9241-210. Este trabalho apresenta o documento formal de modelagem de Personas, Mapas de Empatia e Jornadas do Usuário do sistema. A metodologia empírica fundamenta-se nas teorias de Alan Cooper (\textit{Goal-Directed Design}), Dave Gray e Alexander Osterwalder (\textit{Empathy Map Canvas}), Don Norman (Design Emocional) e Jakob Nielsen (Engenharia de Usabilidade). São especificadas três personas arquetípicas: a persona primária \textit{Beatriz Matos} (Discente de Enfermagem em formação hospitalar), a persona secundária \textit{Prof. Dr. Ricardo Sobral} (Docente e Auditor Sênior de Segurança do Paciente) e a persona terciária \textit{Carlos Eduardo Santos} (Administrador de TI e Suporte Acadêmico), complementadas pela caracterização de uma anti-persona (O Auditor Punitivo Tradicional). Para cada perfil, são construídos mapas de empatia em seis quadrantes e diagramas vetoriais de jornada com curvas emocionais e pontos de contato. A partir dos achados de usabilidade, derivam-se diretrizes ergonômicas de interface clínica, padrões de densidade informacional para estações desktop e quatro matrizes de rastreabilidade bidirecional cruzando dores e ganhos com os vinte e cinco requisitos funcionais (RF01 a RF25), dez requisitos não funcionais (RNF01 a RNF10) e vinte e dois casos de uso (UC01 a UC22) do SIGEA-GTT. O documento é concluído com a avaliação heurística de usabilidade e critérios rigorosos do \textit{Quality Gate} acadêmico.

\textbf{Palavras-chave}: Design Centrado no Usuário. Personas. Mapa de Empatia. Jornada do Usuário. Global Trigger Tool. Ergonomia de Software. Usabilidade Clínica. SIGEA-GTT.
\end{resumo}
